% TOP_PROBLEM: {"id":30007739,"problem_number":"LOCAL-30007739","title":"Origins of child penalties","category_id":21,"category":{"id":21,"name":"economics","display_name":"Economics","description":"Open economic theory statements, published research agendas, and proposed scoped specifications; question provenance and historical priority are qualified per record.","slug":"economics","order_index":21,"created_at":"2026-10-08T00:00:00Z"},"status":"provenance_pending","external_url":null,"legacy_ids":[],"published":false,"statement":"Identify how one childcare or paternal-leave policy changes mothers’ versus fathers’ earnings losses after first birth under an explicit no-birth counterfactual.","economics":{"original_id":228,"field":"Education, families and demography","kind":"identification","formalization_basis":"proposed_specification","source_status":"not_verified","priority_status":"not_established","priority_note":"No readable primary passage explicitly stating this precise open question was verified. A supporting paper, abstract, or topic citation is insufficient to establish origin or unresolved status.","earliest_verified_reference":null,"current_reference":null,"formalization_note":"Proposed scoped research specification. The original catalogue prompt was a synthesis, and no canonical conjecture or verified original formulation is claimed. The supporting source, when retained, establishes context or existing results rather than this exact open question."},"statusQualification":"Provenance pending: an explicit published statement of this exact open question was not verified. This is a catalogue-authored candidate specification, not a certified open conjecture. Proposed scoped research specification. The original catalogue prompt was a synthesis, and no canonical conjecture or verified original formulation is claimed. The supporting source, when retained, establishes context or existing results rather than this exact open question.","statusReviewedAt":"2026-10-08"}
% FIRST_POSTED: 2026-10-08
% ENTRY_KIND: statement_only
% !TeX program = lualatex
\documentclass[11pt]{article}
\usepackage{fontspec}
\usepackage[a4paper,margin=25mm]{geometry}
\usepackage{amsmath,amssymb,mathtools}
\usepackage{xurl}
\usepackage[hidelinks]{hyperref}
\newcommand{\sref}[2]{\hyperlink{source:#1:#2}{[S#2]}}
\setlength{\emergencystretch}{3em}
\begin{document}
% problemId:30007739
\section*{Origins of child penalties}\label{30007739}
\subsection{Definitions and mathematical statement}
Consider couples expecting a first child in one country, with both parents employed before childbirth. Let \(a\in\{0,1\}\) denote a specified policy that changes childcare access or fathers’ nontransferable leave, and let \(Y_{igt}(a)\) be real earnings for parent \(g\) in couple \(i\) at event time \(t\), where \(g\in\{m,f\}\) indexes mothers and fathers and \(t=0\) is the first birth. Let \(\widetilde Y_{igt}(a)\) denote earnings under a specified no-birth counterfactual. Define the gendered parenthood gap at each \(t\in\{1,\ldots,10\}\) by
\[P_t(a)=E[(Y_{imt}(a)-\widetilde Y_{imt}(a))-(Y_{ift}(a)-\widetilde Y_{ift}(a))].\]
Identify \(P_t(1)-P_t(0)\) using credible policy variation and a defensible no-birth trajectory, rather than attribute the entire event-study gap to discrimination or preferences. Collect employer responses, hours, childcare use, partner time, and stated expectations; distinguish mechanisms only with additional variation or explicit exclusion restrictions. Address selection into parenthood and leave eligibility, and retain parents who leave employment. The catalogue sources document family policies and child penalties, but an explicit open passage establishing this particular mechanism decomposition was not verified. This specification asks a scoped identification question and makes no claim that child penalties themselves are unknown.

\par\textbf{Formulation and scope.} Proposed scoped research specification. The original catalogue prompt was a synthesis, and no canonical conjecture or verified original formulation is claimed. The supporting source, when retained, establishes context or existing results rather than this exact open question.

\par\textbf{Open-question provenance.} An explicit published open-question statement was not verified. Sources below are supporting literature and must not be treated as a first listing.

\par\textbf{Historical priority.} No readable primary passage explicitly stating this precise open question was verified. A supporting paper, abstract, or topic citation is insufficient to establish origin or unresolved status.

\subsection{Short English statement}
Identify how one childcare or paternal-leave policy changes mothers’ versus fathers’ earnings losses after first birth under an explicit no-birth counterfactual.

\subsection{Sources}
\begin{itemize}\raggedright
\item[\textnormal{[S1]}]\hypertarget{source:30007739:1}{} Year not verified. Families, Labor Markets, and Policy (2022). Suggested supporting reference from the initial catalogue; exact open-question passage not verified..
\url{https://www.nber.org/papers/w30685}.
{\small\textit{Unverified bibliographic lead Bibliographic lead only. It is not evidence of an original explicit open-question statement.}}

\item[\textnormal{[S2]}]\hypertarget{source:30007739:2}{} Year not verified. The Child Penalty Atlas (2025). Suggested supporting reference from the initial catalogue; exact open-question passage not verified..
\url{https://academic.oup.com/restud/article/92/5/3174/7840285}.
{\small\textit{Unverified bibliographic lead Bibliographic lead only. It is not evidence of an original explicit open-question statement.}}
\end{itemize}
\end{document}
